% GENERATED FILE. Edit canonical lessons, then run npm run generate:books.
% canonical-source-hash: fnv1a64:ff581e069c97cf9b
% canonical-lessons: TA-C524-callatai, TA-C524-ural, TA-C524-ulakkai, TA-C524-kottu, TA-C524-cuttikari

\chapter{Kitchen Work}
\label{ch:ta-kitchen-work}
\hlchaptermodality{524}

\begin{chapteropening}
\textbf{By the end of this chapter:} \emph{I can talk about old kitchen tools and kitchen jobs.}

Complete the last lesson of chapter 524: I can talk about old kitchen tools and kitchen jobs.
\end{chapteropening}

\section[\texorpdfstring{\ta{சல்லடை} (callaṭai)}{சல்லடை (callaṭai)}]{\ta{சல்லடை} (callaṭai) — a sieve}

\label{lesson:TA-C524-callatai}



\begin{quote}
Before the new one: say the Tamil for a drain, then the Tamil for to filter.
\end{quote}

\subsection*{You'll want to know: \ta{சல்லடை}}
\textbf{\ta{சல்லடை}} — \emph{callaṭai} — \textquotedblleft{}a sieve\textquotedblright{}.

\begin{quote}
\ta{பாட்டி} \ta{மாவைச்} \ta{சல்லடையில்} \ta{சலித்தாள்}.
\end{quote}

\begin{quote}
\emph{pāṭṭi māvaic callaṭaiyil calittāḷ.}
\end{quote}

\begin{quote}
Grandmother sifted the flour in a sieve.
\end{quote}

The verb that goes with it is \ta{சலி}, to sift: \ta{மாவைச்} \ta{சலி}.

\begin{cousinweb}
\textbf{\ta{சல்லடை}} (\emph{callaṭai}) is native Dravidian: DEDR 2370 gives \ta{சலி} (\emph{cali}), \textquotedblleft{}to sift\textquotedblright{}, and \ta{சல்லடை} (\emph{callaṭai}), \textquotedblleft{}a sieve\textquotedblright{}.
\end{cousinweb}

\begin{grammarlens}[title={What you've built}]
One more word for describing what you lived and felt.
\end{grammarlens}

\subsection*{Your turn}
\begin{itemize}
\raggedright
  \item \emph{Say it:} \emph{callaṭai}
  \item \emph{Say it:} \emph{callaṭai}, once more
  \item \emph{Say it:} \emph{pāṭṭi māvaic callaṭaiyil calittāḷ.} — Grandmother sifted the flour in a sieve.
  \item \emph{Recall:} say the Tamil for a drain, then the Tamil for to filter, then say \emph{callaṭai} again
\end{itemize}

\begin{tcolorbox}[breakable,colback=teal!4,colframe=teal!35!black,title={Before you move on}]
What does \ta{சல்லடை} mean? (\textquotedblleft{}A sieve\textquotedblright{}.) Say it once more.
\end{tcolorbox}
\section[\texorpdfstring{\ta{உரல்} (ural)}{உரல் (ural)}]{\ta{உரல்} (ural) — a mortar (for pounding)}

\label{lesson:TA-C524-ural}



\begin{quote}
Before the new one: say the Tamil for to filter, then the Tamil for a sieve.
\end{quote}

\subsection*{You'll want to know: \ta{உரல்}}
\textbf{\ta{உரல்}} — \emph{ural} — \textquotedblleft{}a mortar (for pounding)\textquotedblright{}.

\begin{quote}
\ta{முன்பு} \ta{அரிசியை} \ta{உரலில்} \ta{இடித்தார்கள்}.
\end{quote}

\begin{quote}
\emph{muṉpu ariciyai uralil iṭittārkaḷ.}
\end{quote}

\begin{quote}
In the past they pounded rice in a mortar.
\end{quote}

Grain is pounded in an \ta{உரல்} with a heavy pole.

\begin{cousinweb}
\textbf{\ta{உரல்}} (\emph{ural}) is native Dravidian: DEDR 651 gives \ta{உரல்} (\emph{ural}), \textquotedblleft{}a mortar\textquotedblright{}, with cognates across the south.
\end{cousinweb}

\begin{grammarlens}[title={What you've built}]
One more word for describing what you lived and felt.
\end{grammarlens}

\subsection*{Your turn}
\begin{itemize}
\raggedright
  \item \emph{Say it:} \emph{ural}
  \item \emph{Say it:} \emph{ural}, once more
  \item \emph{Say it:} \emph{muṉpu ariciyai uralil iṭittārkaḷ.} — In the past they pounded rice in a mortar.
  \item \emph{Recall:} say the Tamil for to filter, then the Tamil for a sieve, then say \emph{ural} again
\end{itemize}

\begin{tcolorbox}[breakable,colback=teal!4,colframe=teal!35!black,title={Before you move on}]
What does \ta{உரல்} mean? (\textquotedblleft{}A mortar (for pounding)\textquotedblright{}.) Say it once more.
\end{tcolorbox}
\section[\texorpdfstring{\ta{உலக்கை} (ulakkai)}{உலக்கை (ulakkai)}]{\ta{உலக்கை} (ulakkai) — a pestle}

\label{lesson:TA-C524-ulakkai}



\begin{quote}
Before the new one: say the Tamil for a sieve, then the Tamil for a mortar.
\end{quote}

\subsection*{You'll want to know: \ta{உலக்கை}}
\textbf{\ta{உலக்கை}} — \emph{ulakkai} — \textquotedblleft{}a pestle\textquotedblright{}.

\begin{quote}
\ta{அவள்} \ta{உலக்கையால்} \ta{நெல்லைக்} \ta{குத்தினாள்}.
\end{quote}

\begin{quote}
\emph{avaḷ ulakkaiyāl nellaik kuttiṉāḷ.}
\end{quote}

\begin{quote}
She pounded the paddy with a pestle.
\end{quote}

The tool takes -\ta{ஆல்}, \textquotedblleft{}with\textquotedblright{}: \ta{உலக்கையால்} \ta{குத்து}, to pound with a pestle.

\begin{cousinweb}
\textbf{\ta{உலக்கை}} (\emph{ulakkai}) is native Dravidian: DEDR 672 gives \ta{உலக்கை} (\emph{ulakkai}), \textquotedblleft{}a pestle\textquotedblright{}.
\end{cousinweb}

\begin{grammarlens}[title={What you've built}]
One more word for describing what you lived and felt.
\end{grammarlens}

\subsection*{Your turn}
\begin{itemize}
\raggedright
  \item \emph{Say it:} \emph{ulakkai}
  \item \emph{Say it:} \emph{ulakkai}, once more
  \item \emph{Say it:} \emph{avaḷ ulakkaiyāl nellaik kuttiṉāḷ.} — She pounded the paddy with a pestle.
  \item \emph{Recall:} say the Tamil for a sieve, then the Tamil for a mortar, then say \emph{ulakkai} again
\end{itemize}

\begin{tcolorbox}[breakable,colback=teal!4,colframe=teal!35!black,title={Before you move on}]
What does \ta{உலக்கை} mean? (\textquotedblleft{}A pestle\textquotedblright{}.) Say it once more.
\end{tcolorbox}
\section[\texorpdfstring{\ta{கொட்டு} (koṭṭu)}{கொட்டு (koṭṭu)}]{\ta{கொட்டு} (koṭṭu) — to pour out; to empty out}

\label{lesson:TA-C524-kottu}



\begin{quote}
Before the new one: say the Tamil for a mortar, then the Tamil for a pestle.
\end{quote}

\subsection*{You'll want to know: \ta{கொட்டு}}
\textbf{\ta{கொட்டு}} — \emph{koṭṭu} — \textquotedblleft{}to pour out; to empty out\textquotedblright{}.

\begin{quote}
\ta{அவன்} \ta{பையிலிருந்த} \ta{அரிசியை} \ta{டப்பாவில்} \ta{கொட்டினான்}.
\end{quote}

\begin{quote}
\emph{avaṉ paiyilirunta ariciyai ṭappāvil koṭṭiṉāṉ.}
\end{quote}

\begin{quote}
He emptied the rice from the bag into the tin.
\end{quote}

Past (I …): \ta{கொட்டினேன்}, \emph{koṭṭiṉēṉ}.

Not a form of \ta{கொடு} (\emph{koṭu}), to give. A separate verb of the same shape means \textquotedblleft{}to sting\textquotedblright{} (DEDR 2064): \ta{தேள்} \ta{கொட்டியது}, a scorpion stung.

\begin{cousinweb}
\textbf{\ta{கொட்டு}} (\emph{koṭṭu}) is native Dravidian: DEDR 2065 gives \ta{கொட்டு} (\emph{koṭṭu}), \textquotedblleft{}to pour forth, to empty the contents of a basket or sack\textquotedblright{}.
\end{cousinweb}

\begin{grammarlens}[title={What you've built}]
One more word for describing what you lived and felt.
\end{grammarlens}

\subsection*{Your turn}
\begin{itemize}
\raggedright
  \item \emph{Say it:} \emph{koṭṭu}
  \item \emph{Say it:} \emph{koṭṭu}, once more
  \item \emph{Say it:} \emph{avaṉ paiyilirunta ariciyai ṭappāvil koṭṭiṉāṉ.} — He emptied the rice from the bag into the tin.
  \item \emph{Recall:} say the Tamil for a mortar, then the Tamil for a pestle, then say \emph{koṭṭu} again
\end{itemize}

\begin{tcolorbox}[breakable,colback=teal!4,colframe=teal!35!black,title={Before you move on}]
What does \ta{கொட்டு} mean? (\textquotedblleft{}To pour out; to empty out\textquotedblright{}.) Say it once more.
\end{tcolorbox}
\section[\texorpdfstring{\ta{சுத்திகரி} (cuttikari)}{சுத்திகரி (cuttikari)}]{\ta{சுத்திகரி} (cuttikari) — to purify}

\label{lesson:TA-C524-cuttikari}



\begin{quote}
Before the new one: say the Tamil for a pestle, then the Tamil for to pour out.
\end{quote}

\subsection*{You'll want to know: \ta{சுத்திகரி}}
\textbf{\ta{சுத்திகரி}} — \emph{cuttikari} — \textquotedblleft{}to purify\textquotedblright{}.

\begin{quote}
\ta{இந்த} \ta{இயந்திரம்} \ta{குடிநீரைச்} \ta{சுத்திகரிக்கிறது}.
\end{quote}

\begin{quote}
\emph{inta iyantiram kuṭinīraic cuttikarikkiṟatu.}
\end{quote}

\begin{quote}
This machine purifies drinking water.
\end{quote}

Past (I …): \ta{சுத்திகரித்தேன்}, \emph{cuttikarittēṉ}.

What is purified takes -\ta{ஐ}: \ta{நீரைச்} \ta{சுத்திகரி}.

\begin{cousinweb}
\textbf{\ta{சுத்திகரி}} (\emph{cuttikari}) is built on Sanskrit \emph{śuddhi}, \textquotedblleft{}cleansing, purity\textquotedblright{} (CDIAL 12523), with -\ta{கரி}, \textquotedblleft{}to make\textquotedblright{}, from Sanskrit \emph{kṛ}: \textquotedblleft{}to make pure\textquotedblright{}.
\end{cousinweb}

\begin{grammarlens}[title={What you've built}]
One more word for describing what you lived and felt.
\end{grammarlens}

\subsection*{Your turn}
\begin{itemize}
\raggedright
  \item \emph{Say it:} \emph{cuttikari}
  \item \emph{Say it:} \emph{cuttikari}, once more
  \item \emph{Say it:} \emph{inta iyantiram kuṭinīraic cuttikarikkiṟatu.} — This machine purifies drinking water.
  \item \emph{Recall:} say the Tamil for a pestle, then the Tamil for to pour out, then say \emph{cuttikari} again
\end{itemize}

\begin{tcolorbox}[breakable,colback=teal!4,colframe=teal!35!black,title={Before you move on}]
What does \ta{சுத்திகரி} mean? (\textquotedblleft{}To purify\textquotedblright{}.) Say it once more.
\end{tcolorbox}
